
%\newpage
\section{Frequenzmessung}
\label{sec:frequency}

Beginnend mit Version 1.10k kann die Frequenzmessung mit einem Bedienmenü angewählt werden.
Die normale Frequenzmessung wird durch Zählen der fallenden Flanken des Eingangssignals T0 (PD4)
mit dem Zähler 0 (COUNTER0) für eine Sekunde erledigt. Für die Einhaltung einer exakten Sekunde
wird der Zähler 1 mit einem 256:1 Vorteiler der CPU-Frequenz benutzt. Der 16-Bit Zähler des ATmega
kann mit dem Vorteiler auch bei 16 MHz CPU-Frequenz in einem Durchlauf eine Sekunde auszählen.
Für das Starten und das Stoppen des Zählers 0 werden die Vergleichregister B und A des Zählers 1
benutzt. Damit keine zusätzliche Zeitunsicherheit bei den Abfragen entsteht, werden die
Interrupt-Service-Routinen für beide Vergleichs-Ereignisse benutzt.
Die Zeitverzögerung durch die beiden Interrupt-Service-Routinen ist ungefähr gleich.
Für die Einhaltung der exakten Sekunde ist eine konstante Zeitverzögerung unerheblich.
Durch Analyse des Assemblercodes kann die Zeitungleichheit ausgeglichen werden.\\

Bei Frequenzen unter \(33kHz\) wird die Messung durch die Messung einer Periodendauer
ergänzt. Diese Messung wird im Anschluss an eine normale Frequenzmessung durchgeführt.
Dabei werden die Zeit einer bestimmten Anzahl von Pegelwechsel des PCINT20-Eingangs (PD4) 
mit dem Zähler 0 gemessen. 
Bei der Periodenmessung sollte sowohl die negative Pulsweite als auch die positive Pulsweite
in jeder Periode mindestens \(10\mu s\) betragen.
Der Zähler 0 läuft dabei mit voller CPU-Taktrate und das ergibt für eine Periode eine
Aulösung von \(125ns\). Durch eine größere Anzahl von gemessenen Perioden kann die Auflösung
verbessert werden. Bei 125 Perioden (250 Pegelwechsel) ergeben sich so schon eine mittlere
Auflösung für eine Periode von \(1ns\). Damit keine Ungenauigkeit beim Starten und
Stoppen von Zähler 0 entsteht, wird der Zähler 0 mit dem ersten Pegelwechsel von
PCINT20 gestartet und nach der vorgegebenen Zahl über die gleiche Interrupt Service Routine
wieder gestoppt.
Die Zahl der Perioden wird so gewählt, dass die Messzeit etwa 10 Millionen Takte beträgt.
Der Fehleranteil eines Taktes macht dann 0,1~ppm aus.
Bei \(8MHz\) beträgt die Messzeit also etwa 1,25 Sekunden.
Aus der so ermittelten mittleren Periode wird dann eine Frequenz mit höherer Auflösung berechnet.

Zur Kontrolle wurden zwei Tester gegeneinander gemessen.
Erst wurden mit Tester 2 die Testfrequenzen erzeugt und mit Tester 1 die Frequenzen gemessen.
Danach wurde die Messung mit vertauschten Testern wiederholt.
Abbildung \ref{fig:freq-ppm} zeigt die Ergebnisse.
Die nahezu konstanten relativen Abweichungen sind durch die geringen Frequenzunterschiede der beiden Quarze zu erklären.

\begin{figure}[H]
\centering
% GNUPLOT: LaTeX picture with Postscript
\begingroup
  \makeatletter
  \providecommand\color[2][]{%
    \GenericError{(gnuplot) \space\space\space\@spaces}{%
      Package color not loaded in conjunction with
      terminal option `colourtext'%
    }{See the gnuplot documentation for explanation.%
    }{Either use 'blacktext' in gnuplot or load the package
      color.sty in LaTeX.}%
    \renewcommand\color[2][]{}%
  }%
  \providecommand\includegraphics[2][]{%
    \GenericError{(gnuplot) \space\space\space\@spaces}{%
      Package graphicx or graphics not loaded%
    }{See the gnuplot documentation for explanation.%
    }{The gnuplot epslatex terminal needs graphicx.sty or graphics.sty.}%
    \renewcommand\includegraphics[2][]{}%
  }%
  \providecommand\rotatebox[2]{#2}%
  \@ifundefined{ifGPcolor}{%
    \newif\ifGPcolor
    \GPcolortrue
  }{}%
  \@ifundefined{ifGPblacktext}{%
    \newif\ifGPblacktext
    \GPblacktexttrue
  }{}%
  % define a \g@addto@macro without @ in the name:
  \let\gplgaddtomacro\g@addto@macro
  % define empty templates for all commands taking text:
  \gdef\gplbacktext{}%
  \gdef\gplfronttext{}%
  \makeatother
  \ifGPblacktext
    % no textcolor at all
    \def\colorrgb#1{}%
    \def\colorgray#1{}%
  \else
    % gray or color?
    \ifGPcolor
      \def\colorrgb#1{\color[rgb]{#1}}%
      \def\colorgray#1{\color[gray]{#1}}%
      \expandafter\def\csname LTw\endcsname{\color{white}}%
      \expandafter\def\csname LTb\endcsname{\color{black}}%
      \expandafter\def\csname LTa\endcsname{\color{black}}%
      \expandafter\def\csname LT0\endcsname{\color[rgb]{1,0,0}}%
      \expandafter\def\csname LT1\endcsname{\color[rgb]{0,1,0}}%
      \expandafter\def\csname LT2\endcsname{\color[rgb]{0,0,1}}%
      \expandafter\def\csname LT3\endcsname{\color[rgb]{1,0,1}}%
      \expandafter\def\csname LT4\endcsname{\color[rgb]{0,1,1}}%
      \expandafter\def\csname LT5\endcsname{\color[rgb]{1,1,0}}%
      \expandafter\def\csname LT6\endcsname{\color[rgb]{0,0,0}}%
      \expandafter\def\csname LT7\endcsname{\color[rgb]{1,0.3,0}}%
      \expandafter\def\csname LT8\endcsname{\color[rgb]{0.5,0.5,0.5}}%
    \else
      % gray
      \def\colorrgb#1{\color{black}}%
      \def\colorgray#1{\color[gray]{#1}}%
      \expandafter\def\csname LTw\endcsname{\color{white}}%
      \expandafter\def\csname LTb\endcsname{\color{black}}%
      \expandafter\def\csname LTa\endcsname{\color{black}}%
      \expandafter\def\csname LT0\endcsname{\color{black}}%
      \expandafter\def\csname LT1\endcsname{\color{black}}%
      \expandafter\def\csname LT2\endcsname{\color{black}}%
      \expandafter\def\csname LT3\endcsname{\color{black}}%
      \expandafter\def\csname LT4\endcsname{\color{black}}%
      \expandafter\def\csname LT5\endcsname{\color{black}}%
      \expandafter\def\csname LT6\endcsname{\color{black}}%
      \expandafter\def\csname LT7\endcsname{\color{black}}%
      \expandafter\def\csname LT8\endcsname{\color{black}}%
    \fi
  \fi
  \setlength{\unitlength}{0.0500bp}%
  \begin{picture}(7200.00,5040.00)%
    \gplgaddtomacro\gplbacktext{%
      \csname LTb\endcsname%
      \put(814,704){\makebox(0,0)[r]{\strut{}-25}}%
      \csname LTb\endcsname%
      \put(814,1156){\makebox(0,0)[r]{\strut{}-20}}%
      \csname LTb\endcsname%
      \put(814,1609){\makebox(0,0)[r]{\strut{}-15}}%
      \csname LTb\endcsname%
      \put(814,2061){\makebox(0,0)[r]{\strut{}-10}}%
      \csname LTb\endcsname%
      \put(814,2513){\makebox(0,0)[r]{\strut{}-5}}%
      \csname LTb\endcsname%
      \put(814,2966){\makebox(0,0)[r]{\strut{} 0}}%
      \csname LTb\endcsname%
      \put(814,3418){\makebox(0,0)[r]{\strut{} 5}}%
      \csname LTb\endcsname%
      \put(814,3870){\makebox(0,0)[r]{\strut{} 10}}%
      \csname LTb\endcsname%
      \put(814,4323){\makebox(0,0)[r]{\strut{} 15}}%
      \csname LTb\endcsname%
      \put(814,4775){\makebox(0,0)[r]{\strut{} 20}}%
      \csname LTb\endcsname%
      \put(946,484){\makebox(0,0){\strut{} 1}}%
      \csname LTb\endcsname%
      \put(1783,484){\makebox(0,0){\strut{} 10}}%
      \csname LTb\endcsname%
      \put(2619,484){\makebox(0,0){\strut{} 100}}%
      \csname LTb\endcsname%
      \put(3456,484){\makebox(0,0){\strut{} 1000}}%
      \csname LTb\endcsname%
      \put(4293,484){\makebox(0,0){\strut{} 10000}}%
      \csname LTb\endcsname%
      \put(5130,484){\makebox(0,0){\strut{} 100000}}%
      \csname LTb\endcsname%
      \put(5966,484){\makebox(0,0){\strut{} 1e+06}}%
      \csname LTb\endcsname%
      \put(6803,484){\makebox(0,0){\strut{} 1e+07}}%
      \put(176,2739){\rotatebox{-270}{\makebox(0,0){\strut{}Error / ppm}}}%
      \put(3874,154){\makebox(0,0){\strut{}frequency / Hz}}%
    }%
    \gplgaddtomacro\gplfronttext{%
    }%
    \gplbacktext
    \put(0,0){\includegraphics{../GNU/frequency-ppm}}%
    \gplfronttext
  \end{picture}%
\endgroup

\caption{Relativer Fehler für Frequenzmessung }
\label{fig:freq-ppm}
\end{figure}

\subsection{Kalibration der Frequenz mit GPS- oder GLONASS-Empfänger}

Eine Abstimmung der Quarzfrequenz des Transistortesters ist mit einem einstellbaren Kondensator (\(5-25pF\) Trimmer) am Quarz möglich.
Als Referenz zum Abstimmen habe ich das Sekundensignal (1PPS) mit dem GPS-Empfänger {\bf UP501} von {\bf Fastrax Ltd.} und
mit dem GPS/GLONASS-Empfänger {\bf GNS701} von {\bf Global Navigation Systems GmbH} erfolgreich getestet.
Die gemessene Periode ließ sich mit dem Trimmer auf exakt \(1000,000ms\) abgleichen.
Lediglich die letzte angezeigte Ziffer kann um 1 schwanken.
Natürlich ist die Frequenz des Quarzes temperaturabhängig.
Deswegen kann man keine sehr gute Langzeitstabilität erwarten.

Die Abbildung \ref{fig:GPS-1PPS} zeigt die verwendeten Schaltungen mit einem
UM232 USB-seriell-Wandler als Verbindung der Empfänger-Module mit einem Computer.
Der UM232-Wandler versorgt die Schaltung sowohl mit \(5V\) als auch mit \(3,3V\) Betriebsspannung
aus der USB-Versorgung.
Für den Betrieb der Empfänger ist eine Verbindung zum Rechner nicht erforderlich.
Lediglich die \(5V\)-Versorgung der USB-Steckbuchse muss gewährleitet sein.

\begin{figure}[H]
  \begin{subfigure}[b]{9cm}
    \centering
    \includegraphics[width=8.5cm]{../FIG/GPS_UP501.eps}
    \caption{GPS}
  \end{subfigure}
  ~
  \begin{subfigure}[b]{9cm}
    \centering
    \includegraphics[width=8.5cm]{../FIG/GPS_GNS701.eps}
    \caption{GPS/GLONASS}
  \end{subfigure}
  \caption{Sekundenpuls Erzeugung mit GPS-Empfängern }
  \label{fig:GPS-1PPS}
\end{figure}

\subsection{Kalibration der Quarzfrequenz mit einem Uhrenmodul}

Zum Abstimmen der Quarzfrequenz des Transistortesters ist ein Austausch eines der beiden Quarz-Kondensatoren 
gegen einen Trimmer mit einstellbarer Kapazität erforderlich.
Der Vorteil bei der Verwendung von Uhrenmodulen gegenüber von GPS- oder GLONASS-Modulen für den Frequenzabgleich ist,
daß man keine freie Sicht zum Himmel braucht. Man kann den Abgleich nahezu an jedem Ort durchführen.
Untersucht habe ich Uhrenmodule mit einem DS3231 und der Platinen-Aufschrift ,,ZS-042''.
Die untersuchten Module stammen mit großer Wahrscheinlichkeit aus chinesischer Produktion und verwenden
nicht den DS3231SN Chip, sondern einen DS3231M Chip mit einem MEMS Resonator (MEMS = Mico Electro Mechanical System). 
Einen ähnlichen MEMS Resonator benutzt übrigens auch der DCP1301 Chip.
Das Bild~\ref{fig:DS3231M} zeigt eins der verwendeten Module.

\begin{figure}[H]
\centering
\includegraphics[width=9cm]{../PNG/DS3231M.jpg}
\caption{Ein Exemplar der untersuchten DS3231 Module}
\label{fig:DS3231M}
\end{figure}

Der DS3231SN Chip benutzt aber einen Uhrenquarz mit \(32768Hz\) als Zeitbasis.
Bei beiden DS3231 Chip Varianten wird das Temperatur-Verhalten der Schwingfrequenz über eine Chip interne Temperaturmessung
kompensiert, so daß eine gute Frequenzgenauigkeit in einem weiten Temperaturbereich erreicht wird.
Leider ist aber das auf dem Modul herausgeführte \(32kHz\) Signal bei der DS3231M Chip Variante für eine
Kalibration ungeeignet. Bei den vorhandenen vier Modulen waren Frequenzen von \(32641Hz\), \(32710Hz\), \(32730Hz\) und \(32748Hz\) meßbar.
Alle Frequenzen waren also weit entfernt von den eigentlich erwarteten präzisen \(32768Hz\).
Wenn man die Module an einen Arduino UNO anschließt, kann man auch einen 1PPS (1Hz) Ausgang auf den Ausgang SQW schalten.
Dieser Ausgang ist aber so stabil in der Periode, daß er für eine Kalibration benutzt werden kann.
Das Datenblatt des DS3231M vespricht für den 1PPS Ausgang eine Genauigkeit von \(\pm 5ppm\) über den vollen Temperaturbereich
\(-45\celsius\)  bis \(+85\celsius\), während die Genauigkeit des 32kHz Ausgangs nur auf \(\pm 2,5\%\) (\(25000ppm\)) garantiert wird.


Im Datenblatt des DS3231SN wird über einen Temperaturbereich \(-40\celsius\)  bis \(+85\celsius\)  eine Genauigkeit von \(\pm 3.5ppm\)
versprochen und für den Temperaturbereich \(0\celsius\)  bis \(+40\celsius\)  eine Genauigkeit vom \(\pm 2ppm\).
Der DS3231SN Chip verwendet als Zeitbasis einen Uhrenquarz mit der Frequenz \(32768Hz\), der mit schaltbaren Kapazitäten
über einen weiten Temperaturbereich nachgezogen wird. Mit dem bekannten Temperaturgang des Quarzes und einem
auf den Chip integrierten Temperatursensor wird die Frequenz weitgehend konstant gehalten.
Um auch diesen Chip untersuchen zu können, habe ich die DS3231M Chips bei allen vier Modulen gegen DS3231SN Chips
ausgetauscht. Bei einem frisch kalibriertem Transistortester mit \(16MHz\) Quarz ergibt die Frequenzmessung für alle
vier Module die gleiche Frequenz von \(32.76800kHz\) mit einer selten angezeigten Abweichung von \(0.03Hz\).
Das entspricht einer Abweichung von etwa \(1ppm\).
Bruchteile von 1Hz werden übrigens nur angezeigt, wenn die Frequenz aus einer Periodenmessung errechnet wird.
Damit die \(32768Hz\) mit dieser Methode gemessen werden können, habe ich die Grenze für die Periodenmessung
von \(25kHz\) auf \(33kHz\) geändert.
Diese Genauigkeit der Frequenzanzeige kann natürlich nur erreicht werden, wenn die Frequenzquelle entsprechend rauscharm ist.
Ein Signal, welches von einem RC-Generator abgeleitet ist, schwankt erheblich bei den Anzeigewerten.
